% GENERATED FILE. Edit canonical lessons, then run npm run generate:books.
% canonical-source-hash: fnv1a64:8559ba34cae6055b
% canonical-lessons: JA-C468-sakuhin, JA-C468-tenrankai, JA-C468-choukoku, JA-C468-kessaku, JA-C468-kumitateru

\chapter{Works of Art}
\label{ch:ja-works-of-art}
\hlchaptermodality{468}

\begin{chapteropening}
\textbf{By the end of this chapter:} \emph{I can talk about an exhibition.}

Complete the last lesson of chapter 468: I can talk about an exhibition.
\end{chapteropening}

\section[\texorpdfstring{\ja{さくひん} (sakuhin)}{さくひん (sakuhin)}]{\ja{さくひん} (sakuhin) — a work of art; a piece}

\label{lesson:JA-C468-sakuhin}



\begin{quote}
Before the new one: say the Japanese for embroidery, then the Japanese for to whittle.
\end{quote}

\subsection*{You'll want to know: \ja{さくひん}}
\textbf{\ja{さくひん}} — \emph{sakuhin} — \textquotedblleft{}a work of art; a piece\textquotedblright{}.

\begin{quote}
\textbf{\ja{これは} \ja{そふの} \ja{さくひんです。}} — \emph{kore wa sofu no sakuhin desu} — \textquotedblleft{}This is a work by my grandfather.\textquotedblright{}
\end{quote}

\begin{cousinweb}
\textbf{\ja{さくひん}} (\emph{sakuhin}) is Sino-Japanese (kango): two Chinese characters read with their on'yomi, readings borrowed from Chinese. They are \emph{saku}, \textquotedblleft{}make\textquotedblright{} (Middle Chinese \emph{tsak}) and \emph{hin}, \textquotedblleft{}article\textquotedblright{} (\emph{phimX}) (KANJIDIC2; Middle Chinese in Baxter's transcription).
\end{cousinweb}

\begin{grammarlens}[title={What you've built}]
One more word for describing what you lived and felt.
\end{grammarlens}

\subsection*{Your turn}
\begin{itemize}
\raggedright
  \item \emph{Say it:} \emph{sakuhin}
  \item \emph{Say it:} \emph{sakuhin}, once more
  \item \emph{Say it:} \emph{kore wa sofu no sakuhin desu} — This is a work by my grandfather.
  \item \emph{Recall:} say the Japanese for embroidery, then the Japanese for to whittle, then say \emph{sakuhin} again
\end{itemize}

\begin{tcolorbox}[breakable,colback=teal!4,colframe=teal!35!black,title={Before you move on}]
What does \ja{さくひん} mean? (\textquotedblleft{}A work of art; a piece\textquotedblright{}.) Say it once more.
\end{tcolorbox}
\section[\texorpdfstring{\ja{てんらんかい} (tenrankai)}{てんらんかい (tenrankai)}]{\ja{てんらんかい} (tenrankai) — an exhibition}

\label{lesson:JA-C468-tenrankai}



\begin{quote}
Before the new one: say the Japanese for to whittle, then the Japanese for a work of art.
\end{quote}

\subsection*{You'll want to know: \ja{てんらんかい}}
\textbf{\ja{てんらんかい}} — \emph{tenrankai} — \textquotedblleft{}an exhibition\textquotedblright{}.

\begin{quote}
\textbf{\ja{びじゅつかんで} \ja{てんらんかいを} \ja{みました。}} — \emph{bijutsukan de tenrankai o mimashita} — \textquotedblleft{}I saw an exhibition at the art museum.\textquotedblright{}
\end{quote}

\begin{cousinweb}
\textbf{\ja{てんらんかい}} (\emph{tenrankai}) is Sino-Japanese (kango): three Chinese characters read with their on'yomi, readings borrowed from Chinese. They are \emph{ten}, \textquotedblleft{}unfold\textquotedblright{} (Middle Chinese \emph{trjenX}), \emph{ran}, \textquotedblleft{}see\textquotedblright{} (\emph{lamX}), and \emph{kai}, \textquotedblleft{}meeting\textquotedblright{} (\emph{hwajH}) (KANJIDIC2; Middle Chinese in Baxter's transcription).
\end{cousinweb}

\begin{grammarlens}[title={What you've built}]
One more word for describing what you lived and felt.
\end{grammarlens}

\subsection*{Your turn}
\begin{itemize}
\raggedright
  \item \emph{Say it:} \emph{tenrankai}
  \item \emph{Say it:} \emph{tenrankai}, once more
  \item \emph{Say it:} \emph{bijutsukan de tenrankai o mimashita} — I saw an exhibition at the art museum.
  \item \emph{Recall:} say the Japanese for to whittle, then the Japanese for a work of art, then say \emph{tenrankai} again
\end{itemize}

\begin{tcolorbox}[breakable,colback=teal!4,colframe=teal!35!black,title={Before you move on}]
What does \ja{てんらんかい} mean? (\textquotedblleft{}An exhibition\textquotedblright{}.) Say it once more.
\end{tcolorbox}
\section[\texorpdfstring{\ja{ちょうこく} (chōkoku)}{ちょうこく (chōkoku)}]{\ja{ちょうこく} (chōkoku) — a sculpture; carving}

\label{lesson:JA-C468-choukoku}



\begin{quote}
Before the new one: say the Japanese for a work of art, then the Japanese for an exhibition.
\end{quote}

\subsection*{You'll want to know: \ja{ちょうこく}}
\textbf{\ja{ちょうこく}} — \emph{chōkoku} — \textquotedblleft{}a sculpture; carving\textquotedblright{}.

\begin{quote}
\textbf{\ja{こうえんに} \ja{いしの} \ja{ちょうこくが} \ja{あります。}} — \emph{kōen ni ishi no chōkoku ga arimasu} — \textquotedblleft{}There is a stone sculpture in the park.\textquotedblright{}
\end{quote}

\begin{cousinweb}
\textbf{\ja{ちょうこく}} (\emph{chōkoku}) is Sino-Japanese (kango): two Chinese characters read with their on'yomi, readings borrowed from Chinese. They are \emph{chō}, \textquotedblleft{}carve\textquotedblright{} (Middle Chinese \emph{tew}) and \emph{koku}, \textquotedblleft{}time\textquotedblright{} (\emph{khok}) (KANJIDIC2; Middle Chinese in Baxter's transcription).
\end{cousinweb}

\begin{grammarlens}[title={What you've built}]
One more word for describing what you lived and felt.
\end{grammarlens}

\subsection*{Your turn}
\begin{itemize}
\raggedright
  \item \emph{Say it:} \emph{chōkoku}
  \item \emph{Say it:} \emph{chōkoku}, once more
  \item \emph{Say it:} \emph{kōen ni ishi no chōkoku ga arimasu} — There is a stone sculpture in the park.
  \item \emph{Recall:} say the Japanese for a work of art, then the Japanese for an exhibition, then say \emph{chōkoku} again
\end{itemize}

\begin{tcolorbox}[breakable,colback=teal!4,colframe=teal!35!black,title={Before you move on}]
What does \ja{ちょうこく} mean? (\textquotedblleft{}A sculpture; carving\textquotedblright{}.) Say it once more.
\end{tcolorbox}
\section[\texorpdfstring{\ja{けっさく} (kessaku)}{けっさく (kessaku)}]{\ja{けっさく} (kessaku) — a masterpiece}

\label{lesson:JA-C468-kessaku}



\begin{quote}
Before the new one: say the Japanese for an exhibition, then the Japanese for a sculpture.
\end{quote}

\subsection*{You'll want to know: \ja{けっさく}}
\textbf{\ja{けっさく}} — \emph{kessaku} — \textquotedblleft{}a masterpiece\textquotedblright{}.

\begin{quote}
\textbf{\ja{これは} \ja{かれの} \ja{けっさくです。}} — \emph{kore wa kare no kessaku desu} — \textquotedblleft{}This is his masterpiece.\textquotedblright{}
\end{quote}

\begin{cousinweb}
\textbf{\ja{けっさく}} (\emph{kessaku}) is Sino-Japanese (kango): two Chinese characters read with their on'yomi, readings borrowed from Chinese. They are \emph{ke}, \textquotedblleft{}excellence\textquotedblright{} (Middle Chinese \emph{gjet}) and \emph{saku}, \textquotedblleft{}make\textquotedblright{} (\emph{tsak}) (KANJIDIC2; Middle Chinese in Baxter's transcription).
\end{cousinweb}

\begin{grammarlens}[title={What you've built}]
One more word for describing what you lived and felt.
\end{grammarlens}

\subsection*{Your turn}
\begin{itemize}
\raggedright
  \item \emph{Say it:} \emph{kessaku}
  \item \emph{Say it:} \emph{kessaku}, once more
  \item \emph{Say it:} \emph{kore wa kare no kessaku desu} — This is his masterpiece.
  \item \emph{Recall:} say the Japanese for an exhibition, then the Japanese for a sculpture, then say \emph{kessaku} again
\end{itemize}

\begin{tcolorbox}[breakable,colback=teal!4,colframe=teal!35!black,title={Before you move on}]
What does \ja{けっさく} mean? (\textquotedblleft{}A masterpiece\textquotedblright{}.) Say it once more.
\end{tcolorbox}
\section[\texorpdfstring{\ja{くみたてる} (kumitateru)}{くみたてる (kumitateru)}]{\ja{くみたてる} (kumitateru) — to assemble; to put together}

\label{lesson:JA-C468-kumitateru}



\begin{quote}
Before the new one: say the Japanese for a sculpture, then the Japanese for a masterpiece.
\end{quote}

\subsection*{You'll want to know: \ja{くみたてる}}
\textbf{\ja{くみたてる}} — \emph{kumitateru} — \textquotedblleft{}to assemble; to put together\textquotedblright{}.

\begin{quote}
\textbf{\ja{ちちは} \ja{たなを} \ja{くみたてています。}} — \emph{chichi wa tana o kumitatete imasu} — \textquotedblleft{}My father is putting a shelf together.\textquotedblright{}
\end{quote}

Polite \textbf{\ja{くみたてます}}, \emph{kumitatemasu}; an ichidan verb. Past (polite): \textbf{\ja{だいくが} \ja{いすを} \ja{くみたてました。}}, \emph{daiku ga isu o kumitatemashita}.

\begin{cousinweb}
\textbf{\ja{くみたてる}} (\emph{kumitateru}) is a native Japanese compound (wago), inherited rather than borrowed from Chinese: \emph{kumi}, \textquotedblleft{}putting together (kumu)\textquotedblright{} and \emph{tateru}, \textquotedblleft{}to set up\textquotedblright{} (JMdict, KANJIDIC2).
\end{cousinweb}

\begin{grammarlens}[title={What you've built}]
One more word for describing what you lived and felt.
\end{grammarlens}

\subsection*{Your turn}
\begin{itemize}
\raggedright
  \item \emph{Say it:} \emph{kumitateru}
  \item \emph{Say it:} \emph{kumitateru}, once more
  \item \emph{Say it:} \emph{chichi wa tana o kumitatete imasu} — My father is putting a shelf together.
  \item \emph{Recall:} say the Japanese for a sculpture, then the Japanese for a masterpiece, then say \emph{kumitateru} again
\end{itemize}

\begin{tcolorbox}[breakable,colback=teal!4,colframe=teal!35!black,title={Before you move on}]
What does \ja{くみたてる} mean? (\textquotedblleft{}To assemble; to put together\textquotedblright{}.) Say it once more.
\end{tcolorbox}
